\section*{Results}

% ─────────────────────────────────────────────────────────────────────────────
\subsection*{Signal quality and subject inclusion}
% ─────────────────────────────────────────────────────────────────────────────

Of 36 participants in the parent study, 34 met the EMG quality criterion
(median frequency in the physiological range 55--170~Hz and chewing rate
0.6--2.6~Hz in both conditions).
Two participants (M3, M7) were excluded for a confirmed dead electrode in one
condition, rendering their within-subject comparisons invalid
(Supplementary Fig.~\ref{fig:supp_qc}).
Side selection favoured the right masseter in 27 of 34 participants
($N_R$\,=\,27; $N_L$\,=\,7), consistent with the rightward dominance
reported in the parent study~\cite{espinoza2025}.
Full quality metrics per participant are reported in Supplementary Table~S1.

% ─────────────────────────────────────────────────────────────────────────────
\subsection*{Session-averaged null: equivalence across pooled bouts}
% ─────────────────────────────────────────────────────────────────────────────

Averaging across all four chewing bouts per condition yielded no significant
anesthesia-placebo difference for any EMG metric.
Two one-sided equivalence tests (TOST; equivalence bound $d$\,=\,0.5~\cite{Lakens2017})
confirmed formal equivalence for all timing and spectral frequency metrics:
median frequency, mean frequency, chewing rate, cycle duration, and duty cycle
(see Supplementary Table~S3).
Amplitude metrics showed a non-significant trend in the direction of higher values
under anesthesia (median cycle amplitude pooled $d_z$\,$\approx$\,0.15,
$p$\,>\,0.05), replicating the global null for the anesthesia factor reported in
the parent study~\cite{espinoza2025}.
This session-level equivalence is the result that a pooled analysis would
report; the following subsection shows that it reflects dilution of a
transient effect by the pharmacological washout.

% ─────────────────────────────────────────────────────────────────────────────
\subsection*{Bout-resolved analysis: transient amplitude gain at bout~1}
% ─────────────────────────────────────────────────────────────────────────────

Resolving EMG features by chewing bout revealed a pattern consistent with
the washout of topical lidocaine: a substantial effect at bout~1 ($\approx$3~min
post-application) that decayed monotonically across bouts~2--4
(Fig.~\ref{fig:bout_main}).
At bout~1, median cycle amplitude was larger under anesthesia than placebo
($d_z$\,=\,+0.695, Wilcoxon $p$\,=\,0.0003, Benjamini--Hochberg FDR-corrected
$p$\,=\,0.0035~\cite{Benjamini1995}).
The condition\,$\times$\,bout interaction in a repeated-measures ANOVA
confirmed that this difference was not constant across bouts
(Greenhouse--Geisser-corrected interaction $p$\,=\,0.0010~\cite{Pingouin}).
All additional amplitude and power metrics at bout~1 survived FDR correction
(Table~\ref{tab:bout1}): RMS amplitude ($d_z$\,=\,+0.487,
FDR $p$\,=\,0.0132), total power ($d_z$\,=\,+0.456, FDR $p$\,=\,0.0132),
60--150~Hz power ($d_z$\,=\,+0.463, FDR $p$\,=\,0.0189), and 20--60~Hz power
($d_z$\,=\,+0.467, FDR $p$\,=\,0.0274).
Chewing regularity -- indexed as the coefficient of variation of the inter-peak
interval (CV\textsubscript{IPI}) -- was also higher under anesthesia
(lower CV: $d_z$\,=\,$-$0.434, FDR $p$\,=\,0.042), indicating greater
cycle-to-cycle consistency in timing during maximal anesthesia.
By bouts~3--4, all amplitude differences were negligible
($d_z$\,$\leq$\,0.1) and no metric approached significance.
The complete bout-by-bout statistics for all 11 metrics are reported in
Supplementary Table~S4.

\begin{table}[h!]
\centering
\caption{Bout~1 effects of anesthesia vs.\ placebo ($N$\,=\,34).
         FDR-corrected $p$-values computed by the Benjamini--Hochberg
         procedure~\cite{Benjamini1995} across 11 metrics.
         Effect size $d_z$\,=\,paired Cohen's~$d$~\cite{Cohen1988}.
         Direction indicates which condition shows higher values.
         ANE\,=\,anesthesia; PLA\,=\,placebo.}
\label{tab:bout1}
\begin{tabular}{lrrrr}
\toprule
Metric                         & $d_z$   & Raw $p$ & FDR $p$ & Direction     \\
\midrule
Median cycle amplitude         & +0.695  & 0.0003  & 0.0035  & ANE\,$>$\,PLA \\
RMS amplitude                  & +0.487  & 0.0036  & 0.0132  & ANE\,$>$\,PLA \\
Total power                    & +0.456  & 0.0032  & 0.0132  & ANE\,$>$\,PLA \\
Power 60--150~Hz               & +0.463  & 0.0069  & 0.0189  & ANE\,$>$\,PLA \\
Power 20--60~Hz                & +0.467  & 0.0125  & 0.0274  & ANE\,$>$\,PLA \\
CV inter-peak interval         & $-$0.434 & 0.0227 & 0.0415  & ANE\,$<$\,PLA (more regular) \\
\bottomrule
\end{tabular}
\end{table}

% ─────────────────────────────────────────────────────────────────────────────
\subsection*{Washout time course}
% ─────────────────────────────────────────────────────────────────────────────

The anesthesia-placebo difference in median cycle amplitude followed a
monotonic decay across bouts: +30.5, +8.8, $-$4.4, and $-$2.2~\textmu{}V
for bouts~1 through~4, respectively.
The median difference at bout~1 was +22~\textmu{}V (Wilcoxon $p$\,=\,0.0003);
76\% of participants (26/34) showed ANE\,$>$\,PLA at bout~1.
A Spearman trend analysis across the four bout-level estimates confirmed a
significant decreasing trajectory ($p$\,=\,0.0002, slope\,$\approx$\,$-$2.4~\textmu{}V/min).
An exponential decay fitted to the four difference values yielded a time
constant $\tau$\,$\approx$\,2.9~min (illustrative, four points;
Fig.~\ref{fig:washout}A), consistent with topical lidocaine pharmacokinetics
in the oral cavity~\cite{vonGunten2019}.
Robustness of the bout-1 effect across alternative inclusion thresholds is
reported in Supplementary Fig.~\ref{fig:supp_robust}: effect sizes
ranged from $d_z$\,=\,0.65--0.69 for $N$\,=\,30--34 and declined to
$d_z$\,=\,0.26 when both dead-electrode participants were included ($N$\,=\,36),
with the paired Wilcoxon remaining significant ($p$\,=\,0.002),
confirming the QC exclusion rationale.

% ─────────────────────────────────────────────────────────────────────────────
\subsection*{Cycle-locked time-frequency analysis}
% ─────────────────────────────────────────────────────────────────────────────

To characterise the spectrotemporal structure of the amplitude increase at the
single-cycle level, we epoched the bipolar EMG signal around each detected cycle
peak ($[-0.4, +0.4]$~s) and computed Morlet wavelet scalograms across 20--450~Hz.
Grand-average time-frequency maps showed a broadband power burst centred at the
cycle peak that was consistently larger under anesthesia at bout~1 but not at
bout~4 (Fig.~\ref{fig:tf_cluster}).
A cluster-based permutation test~\cite{Maris2007} on the
ANE\,$-$\,PLA difference map at bout~1 identified a single significant cluster
occupying approximately 30--300~Hz from $-$0.15 to +0.1~s around the peak
(a priori directional test, $p$\,=\,0.035; two-tailed $p$\,=\,0.051).
No cluster reached significance at bout~4 ($p_{\min}$\,=\,0.59).
The full 4-bout\,$\times$\,2-condition TF decomposition is provided in
Supplementary Fig.~\ref{fig:supp_tf_full}.

% ─────────────────────────────────────────────────────────────────────────────
\subsection*{Amplitude--timing dissociation}
% ─────────────────────────────────────────────────────────────────────────────

Across all four bouts, no metric indexing masticatory timing differed
significantly between anesthesia and placebo.
At bout~1: median frequency ($d_z$\,=\,$-$0.147, $p$\,=\,0.508), mean
frequency ($d_z$\,=\,$-$0.190, $p$\,=\,0.388), chewing rate
($d_z$\,=\,0.054, $p$\,=\,0.919), and duty cycle ($d_z$\,=\,0.045,
$p$\,=\,0.826) were all non-significant, and all satisfied TOST equivalence
at every bout.
The simultaneous increase in amplitude and power at bout~1 -- alongside
complete stability in frequency and rate at all bouts -- constitutes a
dissociation between the gain and timing dimensions of the masticatory
motor pattern (Fig.~\ref{fig:bout_main}C; Supplementary Table~S4).

% ─────────────────────────────────────────────────────────────────────────────
\subsection*{Exploratory: bilateral masseter synchrony}
% ─────────────────────────────────────────────────────────────────────────────

As an exploratory analysis, bilateral masseter synchrony was examined in
participants with valid signals on both sides ($N$\,=\,31).
Bilateral temporal coupling (envelope Pearson $r$) and amplitude asymmetry
were preserved across conditions and bouts; full results are provided in
Supplementary Fig.~\ref{fig:supp_bilateral}.
In brief, no condition effect was detected at any bout
(all $p_{\text{FDR}}$\,>\,0.59; condition\,$\times$\,bout interaction
$p$\,=\,0.694), indicating that the bilateral temporal structure of the
masticatory pattern was robust to unilateral sensory reduction.

